\chapterbackmatter{Solutions to Supplementary Exercises}

\begin{enumerate}
  \SupplementarySolution{1}{Large Logarithms, Scalar Running, and the Landau Scale}
  \begin{enumerate}
    \item[(a)]
    The relevant expansion parameter is not merely
    \(\lambda/(16\pi^2)\), but
    \[
      x=\frac{\lambda}{16\pi^2}\ln\left(\frac{E^2}{m^2}\right).
    \]
    Fixed-order perturbation theory ceases to be ordered when
    \(|x|\gtrsim1\), even if every loop without its logarithm is small.  At
    \(n\) loops, terms \(x^n\) can then compete with the tree result.  The
    renormalization group replaces the low-scale coupling by a running
    coupling evaluated near \(E\), thereby summing the tower of leading
    logarithms; higher beta-function and anomalous-dimension coefficients
    similarly organize subleading towers.
    \item[(b)]
    Separation of variables gives
    \[
      -\frac{1}{\lambda(\mu)}+\frac{1}{\lambda_0}
        =-\frac{3}{16\pi^2}\ln\left(\frac{\mu}{\mu_0}\right),
    \]
    and hence
    \[
      \boxed{\lambda(\mu)=
        \frac{\lambda_0}{
        1-\dfrac{3\lambda_0}{16\pi^2}
        \ln(\mu/\mu_0)}}.
    \]
    With \(L=\ln(\mu/\mu_0)\), the geometric expansion is
    \[
      \lambda(\mu)=\lambda_0+
        \frac{3\lambda_0^2}{16\pi^2}L+
        \frac{9\lambda_0^3}{(16\pi^2)^2}L^2+O(\lambda_0^4L^3).
    \]
    Thus order \(\lambda_0^{n+1}\) contains the leading logarithm
    \([3L/(16\pi^2)]^n\).
    \item[(c)]
    The denominator of the one-loop solution vanishes at
    \[
      \boxed{\mu_L=\mu_0
        \exp\left(\frac{16\pi^2}{3\lambda_0}\right).}
    \]
    A cutoff must lie below this scale if the one-loop running is to remain
    meaningful.  Conversely, if a cutoff \(\Lambda\) is taken arbitrarily
    high while a finite positive \(\lambda_0\) is held fixed, the trajectory
    encounters the pole.  Avoiding it requires
    \(\lambda_0\lesssim16\pi^2/[3\ln(\Lambda/\mu_0)]\), which tends to zero
    as \(\Lambda\to\infty\).  This is the perturbative form of the
    triviality problem.
  \end{enumerate}

  \SupplementarySolution{2}{Thresholds, Simple Poles, and Characteristic Curves}
  \begin{enumerate}
    \item[(a)]
    One active unit-charge Dirac fermion gives
    \(d\alpha^{-1}/d\ln\mu=-2/(3\pi)\).  There is no running in the
    assumed low-energy effective theory below \(m_1\).  Matching
    continuously gives
    \[
      \alpha^{-1}(m_2)=\alpha^{-1}(\mu_0)
        -\frac{2}{3\pi}\ln\left(\frac{m_2}{m_1}\right).
    \]
    Above \(m_2\), both fermions contribute, so
    \[
      \boxed{\alpha^{-1}(E)=\alpha^{-1}(\mu_0)
        -\frac{2}{3\pi}\left[
          \ln\left(\frac{m_2}{m_1}\right)
          +2\ln\left(\frac{E}{m_2}\right)\right].}
    \]
    Finite matching constants and threshold smoothing are beyond this
    step-function approximation.
    \item[(b)]
    Differentiating at fixed bare coupling gives
    \[
      0=\rho\epsilon\left(g+\frac{ag^2}{\epsilon}\right)
        +\beta(g,\epsilon)
         \left(1+\frac{2ag}{\epsilon}\right)+O(g^3).
    \]
    Put
    \(\beta=-\rho\epsilon g+\beta_2g^2+O(g^3)\).  At order \(g^2\),
    the three contributions are
    \(\rho ag^2+\beta_2g^2-2\rho ag^2\), so
    \(\beta_2=\rho a\).  Therefore
    \[
      \boxed{\beta(g,\epsilon)
        =-\rho\epsilon g+\rho a g^2+O(g^3),\qquad
        \beta(g)=\rho a g^2+O(g^3).}
    \]
    The residue of the simple pole is sufficient at this order.
    \item[(c)]
    Along a characteristic, the boundary coupling evolves according to
    \(dg/dL=bg^2\).  With value \(g\) at \(L=0\), this gives
    \[
      F(L,g)=\frac{g}{1-bgL}.
    \]
    Direct differentiation verifies
    \((-\partial_L+bg^2\partial_g)F=0\) and \(F(0,g)=g\).
    Expanding,
    \[
      \boxed{F(L,g)=\sum_{n=0}^{\infty}b^n g^{n+1}L^n.}
    \]
    Thus the coefficient of \(g^{n+1}L^n\) is \(b^n\).  A one-loop
    beta-function has fixed an infinite sequence of perturbative
    logarithms.
  \end{enumerate}

  \SupplementarySolution{3}{Fixed-Point Stability and Scheme Independence}
  \begin{enumerate}
    \item[(a)]
    Let \(\delta=g-g_*\) and try
    \[
      \delta(\mu)=A\mu^{-\omega}+B\mu^{-2\omega}
        +O(\mu^{-3\omega}).
    \]
    Substitution into
    \(d\delta/d\ln\mu=-\omega\delta+c\delta^2+\cdots\)
    makes the first power automatic.  At order \(\mu^{-2\omega}\),
    \[
      -2\omega B=-\omega B+cA^2,
    \]
    whence
    \[
      \boxed{g(\mu)-g_*
        =A\mu^{-\omega}-\frac{cA^2}{\omega}\mu^{-2\omega}
         +O(\mu^{-3\omega}).}
    \]
    The exponent is fixed by the beta-function slope, while \(A\) labels
    the trajectory.
    \item[(b)]
    The eigenvalues are \(-2\) and \(+1\).  Convenient eigenvectors are
    \[
      V_-=\begin{pmatrix}1\\0\end{pmatrix},\qquad
      V_+=\begin{pmatrix}1\\3\end{pmatrix}.
    \]
    Therefore the displacement from the fixed point is
    \[
      \begin{pmatrix}x(\mu)\\y(\mu)\end{pmatrix}
        =c_-\mu^{-2}V_-+c_+\mu V_+.
    \]
    The positive-eigenvalue component grows toward the ultraviolet.  The
    trajectory reaches the fixed point as \(\mu\to\infty\) precisely when
    \(c_+=0\).  Thus the ultraviolet critical surface is one-dimensional
    in the linear neighborhood and is tangent to \(V_-\).
    \item[(c)]
    The transformed beta-function is
    \[
      \widetilde\beta^\ell(\widetilde{\mathbf g})
        =\frac{\partial f^\ell}{\partial g^m}\beta^m(\mathbf g).
    \]
    At a fixed point the term containing second derivatives of \(f\)
    vanishes because \(\beta^m(\mathbf g_*)=0\).  If
    \(S^\ell{}_m=(\partial f^\ell/\partial g^m)_*\), differentiation gives
    \[
      \widetilde M S=SM,\qquad
      \widetilde M=SMS^{-1}.
    \]
    Similar matrices have identical eigenvalues.  The necessary assumption
    is that the coupling redefinition is differentiable and locally
    nonsingular at the fixed point, so that \(S^{-1}\) exists.
  \end{enumerate}

  \SupplementarySolution{4}{Marginal Directions and Ultraviolet Critical Surfaces}
  \begin{enumerate}
    \item[(a)]
    The solutions are
    \[
      x(\mu)=x_0\frac{\mu_0}{\mu},\qquad
      y(\mu)=\frac{y_0}{
        1-a y_0\ln(\mu/\mu_0)}.
    \]
    The \(x\) direction is ultraviolet attractive.  For \(a>0\), negative
    \(y_0\) approaches zero in the ultraviolet, while positive \(y_0\)
    runs away; the infrared behavior is reversed.  For \(a<0\), the two
    sides interchange.  Since the stability eigenvalue in the \(y\)
    direction is exactly zero, only the quadratic beta-function term
    reveals this marginally attractive or repulsive behavior.
    \item[(b)]
    In the convention
    \(\mu\,d\delta g/d\mu=M\delta g\), a negative eigenvalue gives a
    component that vanishes as \(\mu\to\infty\).  There are three such
    eigenvalues here: \(-4,-1,-0.3\).  Therefore
    \[
      \boxed{\dim\mathcal S_{\rm UV}=3,\qquad
        {\rm codim}\,\mathcal S_{\rm UV}=3.}
    \]
    The coefficients of the three attractive eigenvectors are the three
    independent parameters labeling asymptotically safe trajectories.
    Components along the three positive-eigenvalue directions must be
    tuned to zero.
    \item[(c)]
    With the anomalous dimension convention of the chapter, the
    dimensionless mass satisfies
    \[
      \mu\frac{d\widehat m}{d\mu}
        =(-1-\gamma_m)\widehat m.
    \]
    If \(\gamma_m\) is constant at the fixed point,
    \[
      \boxed{\widehat m(\mu)=\widehat m(\mu_0)
        \left(\frac{\mu}{\mu_0}\right)^{-1-\gamma_m}.}
    \]
    The mass becomes negligible in the ultraviolet when
    \(1+\gamma_m>0\).  If this inequality is reversed it grows as a
    dimensionless perturbation; equality makes it marginal and requires
    examination of higher-order terms.
  \end{enumerate}

  \SupplementarySolution{5}{Banks--Zaks Flow and Short-Distance Forces}
  \begin{enumerate}
    \item[(a)]
    Besides \(g=0\), the two-loop beta-function vanishes at
    \[
      \boxed{\frac{g_*^2}{16\pi^2}=-\frac{b_0}{b_1}.}
    \]
    This is perturbative when \(b_0/|b_1|\ll1\).  Differentiating and using
    the fixed-point relation gives
    \[
      \beta'(g_*)=-3b_0\frac{g_*^2}{16\pi^2}
        -5b_1\frac{g_*^4}{(16\pi^2)^2}
        =\boxed{-\frac{2b_0^2}{b_1}>0}.
    \]
    A displacement obeys
    \(\delta g(\mu)\propto\mu^{\beta'(g_*)}\), so it vanishes as
    \(\mu\to0\).  The Banks--Zaks zero is infrared attractive along the
    gauge-coupling direction.
    \item[(b)]
    In four spacetime dimensions a static potential has mass dimension one.
    At an exact fixed point there is no intrinsic scale, so scale invariance
    permits only
    \[
      \boxed{V(r)=-\frac{\kappa_*}{r},}
    \]
    where the dimensionless coefficient depends on the fixed-point coupling
    and the probe representation.  If an irrelevant displacement satisfies
    \(\delta g(\mu)\propto\mu^\omega\), then at separation \(r\) the natural
    scale is \(\mu\sim r^{-1}\).  Expanding \(\kappa(g)\) therefore gives
    \[
      \boxed{V(r)=-\frac{\kappa_*}{r}
        \left[1+c\,r^{-\omega}
        +O(r^{-2\omega})\right],}
    \]
    with a trajectory-dependent coefficient \(c\).
    \item[(c)]
    The leading momentum-space interaction is proportional to
    \(-C_Fg^2(\mu)/\mathbf q^2\).  Choosing
    \(\mu\sim|\mathbf q|\sim r^{-1}\) moves the large logarithms into the
    running coupling.  For \(b_0>0\),
    \[
      g^2(1/r)=\frac{8\pi^2}{
        b_0\ln[1/(r\Lambda)]}.
    \]
    The Fourier transform of the Coulomb kernel then gives, to leading-log
    accuracy,
    \[
      \boxed{V(r)\simeq-\frac{C_Fg^2(1/r)}{4\pi r}
        =-\frac{2\pi C_F}{
        b_0\,r\ln[1/(r\Lambda)]}.}
    \]
    Corrections from the Fourier transform's range of momenta and from
    higher beta-function coefficients are suppressed by additional inverse
    logarithms.
  \end{enumerate}

  \SupplementarySolution{6}{Heavy-Flavor Matching and Inclusive Running-Coupling Probes}
  \begin{enumerate}
    \item[(a)]
    Write the one-loop solutions on the two sides as
    \[
      \frac{1}{g_H^2(\mu)}
        =\frac{b_H}{8\pi^2}\ln\left(\frac{\mu}{\Lambda_H}\right),
      \qquad
      \frac{1}{g_L^2(\mu)}
        =\frac{b_L}{8\pi^2}\ln\left(\frac{\mu}{\Lambda_L}\right).
    \]
    Continuity at \(\mu=M\) requires
    \(b_H\ln(M/\Lambda_H)=b_L\ln(M/\Lambda_L)\).  Therefore
    \[
      \boxed{\Lambda_L=M
        \left(\frac{\Lambda_H}{M}\right)^{b_H/b_L}.}
    \]
    Here \(b_H=b_0^{(n_f)}\) and
    \(b_L=b_0^{(n_f-1)}\).  Higher orders add a finite matching correction.
    \item[(b)]
    There are two active up-type flavors and three active down-type
    flavors.  Including three colors,
    \[
      R=3\left[2\left(\frac23\right)^2+
        3\left(-\frac13\right)^2\right]
        =\boxed{\frac{11}{3}}.
    \]
    The inclusive sum permits cancellation of soft and collinear
    singularities between real and virtual contributions, so perturbative
    corrections to \(R\) are infrared safe.  A specified isolated quark or
    gluon state is colored and is not an infrared-safe observable; arbitrarily
    soft or collinear radiation changes that exclusive state.
    \item[(c)]
    The electromagnetic current can create a quark--antiquark pair without
    a QCD vertex, so the leading two-parton, and hence two-jet, rate is
    order \(\alpha_s^0\).  A resolved third jet requires emission of a hard
    gluon and therefore one QCD vertex in the amplitude:
    \[
      \frac{\sigma_{3j}(E)}{\sigma_{2j}(E)}
        =C_{\rm jet}\alpha_s(E)+O(\alpha_s^2).
    \]
    For the same infrared-safe jet definition at \(E_1<E_2\), the decrease
    of this ratio tests \(\alpha_s(E_2)<\alpha_s(E_1)\).  No internal
    hadron momentum distribution is required, since that distribution is
    not itself infrared safe.
  \end{enumerate}

  \SupplementarySolution{7}{The Wilson--Fisher Fixed Point and Its Thermal Exponent}
  \begin{enumerate}
    \item[(a)]
    The Gaussian field dimension is
    \([\phi]=(d-2)/2\).  The coupling of \(\phi^{2n}\) consequently has
    infrared RG eigenvalue
    \[
      \boxed{y_{2n}=d-n(d-2).}
    \]
    In \(d=3\), \(y_4=1\), so \(\phi^4\) is relevant, while
    \(y_6=0\), so \(\phi^6\) is marginal by power counting.  In \(d=4\),
    \(y_4=0\), so \(\phi^4\) is marginal, and \(y_6=-2\), so
    \(\phi^6\) is irrelevant.  Nonlinear beta-function terms decide a
    marginal coupling's ultimate behavior.
    \item[(b)]
    The zeros of the beta-function are
    \[
      g_0=0,\qquad
      \boxed{g_*=\frac{\epsilon}{A}+O(\epsilon^2).}
    \]
    Its derivative at the interacting fixed point is
    \[
      \beta'(g_*)=-\epsilon+2Ag_*+O(\epsilon^2)
        =\boxed{\epsilon+O(\epsilon^2)}.
    \]
    With the chapter's momentum-scale convention, this positive eigenvalue
    means that the quartic perturbation decays as the scale is lowered
    toward the infrared fixed point.  It is the leading irrelevant
    direction and controls corrections to scaling.
    \item[(c)]
    The nontrivial zero of the quartic beta-function is
    \[
      g_*=\frac{8\pi^2}{N+8}\epsilon+O(\epsilon^2).
    \]
    Linearization of the mass flow at \(g_2=0\) then gives
    \[
      \boxed{\lambda_t=-2+
        \frac{N+2}{N+8}\epsilon+O(\epsilon^2).}
    \]
    This \(\lambda_t\) is negative in Weinberg's \(\mu\)-flow convention
    and hence is the single relevant infrared perturbation.  In the
    common coarse-graining convention its positive eigenvalue is
    \(y_t=-\lambda_t
      =2-(N+2)\epsilon/(N+8)+O(\epsilon^2)\).
  \end{enumerate}

  \SupplementarySolution{8}{Correlation-Length and Correction-to-Scaling Exponents}
  \begin{enumerate}
    \item[(a)]
    Since \(\nu=1/y_t=-1/\lambda_t\),
    \[
      \boxed{\nu=\frac12+
        \frac{N+2}{4(N+8)}\epsilon+O(\epsilon^2).}
    \]
    Setting \(\epsilon=1\) in this first-order truncation gives
    \[
      \nu_{N=1}\simeq0.5833,\qquad
      \nu_{N=2}\simeq0.6000,\qquad
      \nu_{N=3}\simeq0.6136.
    \]
    These numbers illustrate the trend of the epsilon expansion, but
    omitted higher orders are not parametrically small when
    \(\epsilon=1\), so the displayed precision is not an error estimate.
    \item[(b)]
    For
    \(\beta(g)=-\epsilon g+Ag^2\), the interacting fixed point is
    \(g_*=\epsilon/A\), and
    \[
      \boxed{\omega=\beta'(g_*)=\epsilon+O(\epsilon^2).}
    \]
    Under infrared evolution a displacement in the quartic direction
    decays as \(\mu^\omega\).  Near criticality the inverse correlation
    length is the natural infrared scale, \(\mu\sim\xi^{-1}\).  Therefore a
    regular observable has the form
    \[
      X=X_{\rm leading}
        [1+c\,\xi^{-\omega}+\cdots],
    \]
    with a nonuniversal amplitude \(c\) but a universal exponent \(\omega\).
    \item[(c)]
    At order \(g\), the only one-loop 1PI two-point graph is a tadpole.
    Its loop carries no external momentum, so it changes the mass term but
    cannot generate a divergence proportional to \(p^2\).  Thus
    \(Z_\phi=1+O(g^2)\) and the field anomalous dimension begins at
    \(O(g^2)\).  The first momentum-dependent divergent two-point topology
    is the two-loop sunset graph.  Since \(g_*=O(\epsilon)\),
    \[
      \boxed{\eta=O(g_*^2)=O(\epsilon^2),}
    \]
    so there is no term linear in \(\epsilon\).
  \end{enumerate}

  \SupplementarySolution{9}{Cubic Anisotropy and Competing Fixed Points}
  \begin{enumerate}
    \item[(a)]
    Independent sign flips force every component to occur with even
    exponent, and permutations require symmetric combinations.  At degree
    two the only invariant is \(\phi^2\), which is already \(O(N)\)
    invariant.  At degree four there are two independent invariants,
    \[
      (\phi^2)^2,\qquad \sum_a\phi_a^4.
    \]
    Their difference is a genuine cubic anisotropy.  Higher power sums,
    such as \(\sum_a\phi_a^6\), give further cubic but non-\(O(N)\)
    invariants.
    \item[(b)]
    At the Gaussian point
    \([\phi]=(d-2)/2=1-\epsilon/2\).  A nonderivative monomial of degree
    \(k\) has coupling dimension \(d-k(d-2)/2\).  Apart from the common
    mass deformation, the only cubic-symmetric non-\(O(N)\) operator with
    nonnegative coupling dimension near \(d=4\) is therefore the quartic
    \(\sum_a\phi_a^4\).  Operators of degree six or higher and additional
    derivatives are irrelevant near four dimensions.
    \item[(c)]
    The one-loop four-point fish graphs can be organized by the general
    quartic-coupling tensor formula.  It is convenient to define
    \(\bar u=u/(8\pi^2)\) and \(\bar v=v/(8\pi^2)\), and then suppress the
    bars.  Minimal subtraction gives
    \[
      \boxed{\begin{aligned}
        \beta_u&=-\epsilon u+\frac{N+8}{6}u^2+uv,\\
        \beta_v&=-\epsilon v+2uv+\frac32v^2.
      \end{aligned}}
    \]
    Restoring the unscaled couplings, the quadratic numerators over
    \(48\pi^2\) are
    \((N+8)u^2+6uv\) and \(12uv+9v^2\), respectively.  The absence of
    \(v^2\) in \(\beta_u\) and of \(u^2\) in \(\beta_v\) follows because
    two purely diagonal vertices remain diagonal, whereas two
    \(O(N)\)-invariant vertices cannot generate anisotropy.
    \item[(d)]
    In the rescaled variables the four fixed points are
    \[
      \begin{aligned}
        (u_*,v_*)_{\rm G}&=(0,0),\\
        (u_*,v_*)_{O(N)}&=\left(\frac{6\epsilon}{N+8},0\right),\\
        (u_*,v_*)_{\rm I^N}&=\left(0,\frac{2\epsilon}{3}\right),\\
        (u_*,v_*)_{\rm cub}&=\left(\frac{2\epsilon}{N},
          \frac{2\epsilon(N-4)}{3N}\right).
      \end{aligned}
    \]
    When \(u=0\), the action is a sum of \(N\) independent one-component
    \(\phi^4\) theories.  Its interacting fixed point is therefore \(N\)
    decoupled copies of the Ising Wilson--Fisher theory.
    \item[(e)]
    For \(M_{ij}=\partial\beta_i/\partial g_j|_*\), positive eigenvalues
    decay toward the infrared in the convention
    \(\beta=\mu\,d/d\mu\).  The eigenvalues are
    \[
      \begin{aligned}
        \operatorname{spec}M_{\rm G}&=(-\epsilon,-\epsilon),\\
        \operatorname{spec}M_{O(N)}&=
          \left(\epsilon,\epsilon\frac{4-N}{N+8}\right),\\
        \operatorname{spec}M_{\rm I^N}&=
          \left(\epsilon,-\frac{\epsilon}{3}\right),\\
        \operatorname{spec}M_{\rm cub}&=
          \left(\epsilon,\epsilon\frac{N-4}{3N}\right).
      \end{aligned}
    \]
    Hence the \(O(N)\) point is infrared stable for \(N<4\), while the
    cubic point is stable for \(N>4\).  They merge at \(N=4\), where the
    anisotropy is marginal at this order.  The Gaussian and
    decoupled-Ising points are unstable or saddle points.  A flow sketch
    therefore has the generic critical trajectories ending at \(O(N)\)
    below four components and at the mixed cubic point above four
    components, separated from runaway regions by the stable manifolds of
    the saddles.
    \item[(f)]
    Once a trajectory leaves the stability cone of the quartic form, the
    quartic truncation cannot decide the global potential.  Positive
    sextic terms such as
    \((\phi^2)^3\), \((\phi^2)\sum_a\phi_a^4\), and
    \(\sum_a\phi_a^6\), as well as the bounded microscopic spin length,
    stabilize large fields.  A flow that runs away before reaching a
    scale-invariant fixed point is the perturbative signal of a
    fluctuation-induced first-order transition.  It is therefore a real
    change from the continuous mean-field prediction, not merely a shift
    of critical exponents.
  \end{enumerate}

  \SupplementarySolution{10}{The \(O(N)\) Model at Large \(N\)}
  \begin{enumerate}
    \item[(a)]
    The free kinetic term gives
    \([\phi]=(D-2)/2\).  Since the action is dimensionless,
    \[
      \boxed{[g]=D-4[\phi]=4-D.}
    \]
    Thus the quartic coupling is relevant below four dimensions,
    irrelevant above four dimensions, and classically marginal at \(D=4\).
    \item[(b)]
    Every quartic vertex contributes \(g/N\), whereas every freely summed
    closed \(O(N)\) index loop contributes \(N\).  The leading connected
    vacuum graphs maximize the number of index loops and scale as
    \(N^{L_{\rm index}-V}=N\).  Choosing a coupling independent of \(N\)
    would suppress interactions; omitting the \(1/N\) would make them
    overwhelm the free \(O(N)\) contribution.  The normalization \(g/N\)
    therefore produces a nontrivial saddle and an extensive
    \(O(N)\) free energy.
    \item[(c)]
    Introduce fields \(\rho\) and \(\lambda\) so that
    \(\phi^2/N=\rho\) is enforced by \(\lambda\).  Up to a choice of
    contour for \(\lambda\), the action becomes
    \[
      S_E=\frac12\int d^Dx\,\phi_a(-\partial^2+\lambda)\phi_a
      +N\int d^Dx\left[
        \frac g4\rho^2+\frac{r-\lambda}{2}\rho\right].
    \]
    Integrating out the \(N\) Gaussian fields gives
    \[
      \frac{S_{\rm eff}}N=
      \frac12\operatorname{Tr}\ln(-\partial^2+\lambda)
      +\int d^Dx\left[
        \frac g4\rho^2+\frac{r-\lambda}{2}\rho\right].
    \]
    The overall \(N\) makes the saddle exact at leading order.

    For a translation-invariant symmetric saddle put
    \(m^2=\lambda\).  With a sharp ultraviolet cutoff,
    \[
      \rho=I_D(m),\qquad m^2=r+g\rho,\qquad
      I_D(m)=\int_{|k|<\Lambda}\frac{d^Dk}{(2\pi)^D}
        \frac1{k^2+m^2}.
    \]
    At this order the \(\phi_a\) are free massive fields with propagator
    \(\delta_{ab}/(k^2+m^2)\), while fluctuations of the auxiliary fields
    begin at relative order \(1/N\).

    To display the broken phase, allow
    \(\langle\phi_1\rangle=\sqrt N\,\varphi\).  The saddle equations become
    \[
      m^2\varphi=0,\qquad
      m^2=r+g\bigl[\varphi^2+I_D(m)\bigr].
    \]
    For \(D>2\), define the cutoff-dependent critical parameter
    \[
      r_c=-gI_D(0).
    \]
    If \(r>r_c\), then \(\varphi=0\) and \(m>0\).  At \(r=r_c\), \(m=0\).
    If \(r<r_c\), then \(m=0\) and
    \(\varphi^2=(r_c-r)/g\); the \(N-1\) transverse modes are Goldstone
    bosons.

    Write \(t=r-r_c\).  On the symmetric side,
    \[
      t=m^2+g\,[I_D(0)-I_D(m)].
    \]
    The small-\(m\) behavior separates the requested regimes:
    \[
    \begin{array}{c|c}
      D>4&
      I_D(0)-I_D(m)=c_D\Lambda^{D-4}m^2+\cdots\\
      D=4&
      I_4(0)-I_4(m)=
        \dfrac{m^2}{16\pi^2}\ln\dfrac{\Lambda^2}{m^2}
        +O(m^2)\\[6pt]
      2<D<4&
      I_D(0)-I_D(m)=
        A_Dm^{D-2}+o(m^{D-2})
    \end{array}
    \]
    with
    \[
      A_D=-\frac{\Gamma(1-D/2)}{(4\pi)^{D/2}}>0 .
    \]
    Thus \(D>4\) has mean-field scaling; \(D=4\) has mean-field powers
    multiplied by logarithms; and \(2<D<4\) has interacting large-\(N\)
    scaling.  For \(D\le2\), \(I_D(0)\) is infrared divergent.  No finite
    \(r_c\) supports a massless broken saddle, in accord with the
    lower-critical-dimension obstruction.
    \item[(d)]
    Since \(\xi=m^{-1}\), the preceding gap equation gives
    \[
      \boxed{\nu=
      \begin{cases}
        \dfrac1{D-2},&2<D<4,\\[4pt]
        \dfrac12,&D>4,
      \end{cases}}
    \]
    with \(\nu=1/2\) and logarithmic corrections at \(D=4\).
    At \(D=4-\epsilon\),
    \[
      \nu=\frac1{2-\epsilon}
        =\frac12+\frac{\epsilon}{4}+O(\epsilon^2).
    \]
    The one-loop Wilson--Fisher result is
    \[
      \nu=\frac12+
        \frac{N+2}{4(N+8)}\epsilon+O(\epsilon^2),
    \]
    whose \(N\to\infty\) limit agrees exactly with the expansion above.
  \end{enumerate}

  \SupplementarySolution{11}{The \(\phi^6\) Deformation at Wilson--Fisher}
  Setting \(n=6\) gives
  \[
    \Delta_{\phi^6}=6-3\epsilon+5\epsilon+O(\epsilon^2)
      =6+2\epsilon+O(\epsilon^2).
  \]
  Hence
  \[
    y_6=d-\Delta_{\phi^6}
      =(4-\epsilon)-(6+2\epsilon)+O(\epsilon^2)
      =-2-3\epsilon+O(\epsilon^2).
  \]
  The deformation is irrelevant.  At the Gaussian point its engineering
  eigenvalue is \(d-6(d-2)/2=-2+2\epsilon\), so the interaction makes it
  still more irrelevant near four dimensions.

  This conclusion is local to the ordinary Wilson--Fisher fixed point and
  controlled by the \(4-\epsilon\) expansion.  In \(d=3\), \(\phi^6\) is
  classically marginal; after tuning the mass and quartic coupling, another
  tricritical fixed point can govern the flow.  Irrelevance at one fixed
  point does not rule out that distinct critical surface.

  \SupplementarySolution{12}{A Particle Potential and Weyl Asymptotics}
  \begin{enumerate}
    \item[(a)] With \(\varphi=\sqrt m\,q\),
      \[
        S_E=\int dt\left[\frac12\dot\varphi^2+
          V(\varphi/\sqrt m)\right].
      \]
      Under \(t=b\,t'\), invariance of the kinetic term requires
      \(\varphi(t)=b^{1/2}\varphi'(t')\).  Expanding a smooth potential,
      \(V=\sum_{\mathbf n}g_{\mathbf n}\varphi^{\mathbf n}\), gives
      \[
        g_{\mathbf n}'=b^{\,1+|\mathbf n|/2}g_{\mathbf n}.
      \]
      Every Taylor coefficient therefore grows toward long times and is
      relevant at the free-particle fixed point.  For
      \(V\sim|\varphi|^{-\alpha}\), the coupling scales as
      \(b^{1-\alpha/2}\): it is relevant for \(\alpha<2\), marginal at
      \(\alpha=2\), and irrelevant by this power count for \(\alpha>2\).
    \item[(b)] High energies correspond to short proper time.  The free
      heat trace in a box is
      \[
        Z_0(\beta)
          =L^d\int\frac{d^dp}{(2\pi)^d}
             e^{-\beta p^2/(2m)}
          =L^d\left(\frac{m}{2\pi\beta}\right)^{d/2}.
      \]
      A smooth \(V\) enters the heat-kernel expansion at lower powers of
      \(E\), so the leading coefficient is free.  Since
      \[
        Z(\beta)=\int_0^\infty e^{-\beta E}\,dN(E),
      \]
      a Tauberian inversion of \(Z\sim
      L^d(m/2\pi\beta)^{d/2}\) yields
      \[
        N(E)\sim \frac{L^d}{(2\pi)^d}
          \operatorname{Vol}(B_d)(2mE)^{d/2}.
      \]
      Equivalently, the imaginary part of
      \(\operatorname{Tr}(\omega+i0-H)^{-1}\) gives the density of states;
      representing the resolvent by a proper-time integral leads to the
      same short-time asymptotics.
  \end{enumerate}

\end{enumerate}
